\documentclass[11pt]{article}
\usepackage{amsmath,amssymb,amsthm,booktabs,array,geometry,hyperref,longtable,xcolor,multirow}
\geometry{margin=1in}

\newtheorem{theorem}{Theorem}
\newtheorem{proposition}[theorem]{Proposition}
\newtheorem{definition}[theorem]{Definition}
\newtheorem{remark}[theorem]{Remark}
\newtheorem{corollary}[theorem]{Corollary}
\newtheorem{finding}[theorem]{Finding}

\definecolor{correct}{RGB}{20,130,60}
\definecolor{flagged}{RGB}{180,30,30}
\definecolor{warn}{RGB}{180,110,0}

\newcommand{\ok}[1]{\textcolor{correct}{#1}}
\newcommand{\bad}[1]{\textcolor{flagged}{#1}}
\newcommand{\warn}[1]{\textcolor{warn}{#1}}

\title{SuperBEST Re-Audit: Incorporating Verified Tight Lower Bounds\\
\large Exploration Note --- Not a Table Update}
\author{Monogate Research}
\date{2026-04-22}

\begin{document}
\maketitle

\begin{abstract}
This exploration note re-audits the SuperBEST router using two newly locked lower bounds:
$\mathrm{SB}(\mathrm{mul},\mathrm{general}) = 3$ (Python-certified exhaustive search)
and $\mathrm{SB}(\mathrm{div},\mathrm{general}) \geq 3$ (Lean-verified, \texttt{DivLowerBound3Full.lean}, 0 sorries).
We also correct the positive-domain division cost, which the v5.2 table claimed as $1n$ but
the Lean construction \texttt{div\_two\_node\_pos\_domain} establishes as $2n$.
We do \emph{not} update the canonical SuperBEST table here; all corrections are exploration-level.
The positive-domain $15n$ total is \ok{preserved} under the corrected accounting.
The general-domain total requires upward revision.
\end{abstract}

\tableofcontents

\section{New Locked Lower Bounds (as of 2026-04-22)}

\begin{center}
\begin{tabular}{llll}
\toprule
Bound & Status & Source & Circuit count \\
\midrule
$\mathrm{SB}(\mathrm{mul},\text{general}) = 3$ & \warn{Python-certified} & \texttt{mul\_gen\_tight\_2node\_search.py} & 4112 circuits \\
$\mathrm{SB}(\mathrm{div},\text{general}) \geq 3$ & \ok{Lean-proved} & \texttt{DivLowerBound3Full.lean} & 3072 lemmas \\
$\mathrm{SB}(\mathrm{div},\text{positive}) = 2$ & \ok{Lean-proved (UB)} & \texttt{div\_two\_node\_pos\_domain} & 1 construction \\
\bottomrule
\end{tabular}
\end{center}

\noindent The Lean proofs use the exact F16 operator definitions from \texttt{DivLowerBound.lean}:
\begin{align*}
D_{F1}(x,y) &= \exp(x) - \ln(y) & D_{F9}(x,y) &= x - \ln(y) \\
D_{F10}(x,y) &= x - \ln(-y) & D_{F11}(x,y) &= \ln(\exp(x)+y) \\
D_{F12}(x,y) &= \ln(\exp(x)-y) & D_{F13}(x,y) &= \exp(x \ln y) \\
D_{F14}(x,y) &= \exp(x + \ln y) & D_{F16}(x,y) &= \exp(\ln x + \ln y)
\end{align*}
(plus sign variants F2--F8, F15). Key lemma used throughout: $\ln(-x) = \ln(x)$ for $x \neq 0$
(\texttt{Real.log\_neg\_eq\_log}, Lean 4 Mathlib).

\section{What Changed vs.\ v5.2}

\subsection{General-Domain Multiplication: $2n \to 3n$}

The v5.2 table stated $\mathrm{SB}(\mathrm{mul},\text{general}) \geq 2$ (Lean lower bound).
The Python exhaustive search confirms the exact cost is $3n$:
no 2-node F16 circuit computes $xy$ for all real $x,y$.

\begin{finding}
Every expression involving general-domain multiplication of two free real variables
requires $+1n$ relative to the v5.2 claim.
\end{finding}

\noindent The 3-node upper bound: negate-multiply-negate route gives $xy$ at cost $\leq 5n$,
but the Python search found 3-node constructions. The exact 3-node witness should be
documented in a future Lean proof.

\subsection{General-Domain Division: $2n \to \geq 3n$ (now Lean-proved)}

The v5.2 table stated $\mathrm{SB}(\mathrm{div},\text{general}) \geq 2$ (Lean).
Now \texttt{DivLowerBound3Full.lean} proves $\geq 3$ with 3072 lemmas (0 sorries).
Every 2-node F16 circuit with F1--F12 outer is refuted at $(x,y) = (6,3)$.
Combined with the exp-outer impossibility (F13--F16 always positive, but $x/y$ can be negative),
this completes the $\geq 3$ bound for all outer operators.

\begin{finding}
Every expression involving general-domain division of two free real variables
requires $\geq +1n$ relative to the v5.2 claim.
\end{finding}

\subsection{Positive-Domain Division: Claimed $1n$, Correct $2n$}

The v5.2 table (SuperBEST\_v5.2\_Verified\_Lock.tex) listed:
\begin{center}
\texttt{div($x,y$) \quad 1n (positive) \quad $\geq 2n$ (general)}
\end{center}

\noindent This $1n$ claim is \bad{inconsistent} with the F16 operator definitions.
No single D\_F1 through D\_F16 operator applied to raw inputs $(x,y)$ computes $x/y$ for all $x,y > 0$.
The correct positive-domain cost is $2n$, demonstrated by \texttt{div\_two\_node\_pos\_domain}:
\[
D_{F16}(x,\; D_{F13}(-1,y)) \;=\; \exp(\ln x + \ln(1/y)) \;=\; x/y \quad (x,y > 0).
\]

\begin{finding}
$\mathrm{SB}(\mathrm{div},\text{positive}) = 2$ (not $1$).
The $15n$ positive-domain total is \ok{preserved}: it was computed with $\mathrm{div} = 2n$.
The v5.2 table entry ``$1n$ positive'' for div was an error or notational artifact.
\end{finding}

\noindent\textbf{Why does $15n$ still hold?}
Summing the core 10 operations on positive domain with corrected div:
\[
\underbrace{1}_{\exp} + \underbrace{1}_{\ln} + \underbrace{1}_{\sqrt{\cdot}} + \underbrace{1}_{\mathrm{pow}} + \underbrace{1}_{\mathrm{mul}} + \underbrace{2}_{\mathrm{div}} + \underbrace{2}_{\mathrm{neg}} + \underbrace{2}_{\mathrm{add}} + \underbrace{2}_{\mathrm{sub}} + \underbrace{2}_{\mathrm{abs}} = 15n. \checkmark
\]

\subsection{Positive-Domain Multiplication: Still $1n$}
\[
D_{F16}(x,y) = \exp(\ln x + \ln y) = xy \quad (x,y > 0). \quad\text{(Lean-verified.)}
\]
No change here.

\section{Revised General-Domain Core Table}

\begin{remark}
This table is \emph{exploratory only}. Bounds marked $\geq$ are one-sided.
The canonical table is not updated until all bounds are fully Lean-verified.
\end{remark}

\begin{center}
\begin{tabular}{lrrrll}
\toprule
Operation & Pos.\ (old) & Pos.\ (new) & Gen.\ (old) & Gen.\ (new) & Change? \\
\midrule
$\exp(x)$      & $1n$ & $1n$ & $1n$      & $1n$      & --- \\
$\ln(x)$       & $1n$ & $1n$ & $1n$      & $1n$      & --- \\
$\sqrt{x}$     & $1n$ & $1n$ & ---       & ---       & --- \\
$x^n$          & $1n$ & $1n$ & ---       & ---       & --- \\
$x \cdot y$    & $1n$ & $1n$ & $\geq 2n$ & \bad{$3n$} & \bad{+1n gen.} \\
$x / y$        & \warn{$1n$} & \ok{$2n$} & $\geq 2n$ & \bad{$\geq 3n$} & \bad{+1n pos., +1n gen.} \\
$-x$           & $2n$ & $2n$ & $2n$      & $2n$      & --- \\
$x + y$        & $2n$ & $2n$ & $2n$      & $2n$      & --- \\
$x - y$        & $2n$ & $2n$ & $2n$      & $2n$      & --- \\
$|x|$          & $2n$ & $2n$ & $2n$      & $2n$      & --- \\
\midrule
\textbf{Total (positive)} & $15n$ & $\mathbf{15n}$ & & & \ok{preserved} \\
\textbf{Total (general)}  & & & $\leq 22n$ & $\leq 22n$ & recount needed \\
\bottomrule
\end{tabular}
\end{center}

\noindent The general-domain total with corrected mul=3n and div$\geq$3n:
$1+1+3+3+2+2+2+2 = 16n$ for the 8 operations defined on all of $\mathbb{R}$.
The old ``22n'' figure likely included additional operations or overestimates;
a precise recount is deferred to Core10\_ReCheck.tex.

\section{Impact on Catalog Expressions}

\subsection{Rule: When Does the Correction Apply?}

The $+1n$ corrections for mul and div apply \emph{only} when:
\begin{enumerate}
  \item Both arguments are free variables (not fixed constants), \emph{and}
  \item The domain is not restricted to $(0,\infty) \times (0,\infty)$ for mul.
\end{enumerate}

\noindent\textbf{Unaffected:} $c \cdot x$ with fixed constant $c$ folds into EPL: $D_{F13}(c,x) = x^c$ or $D_{F14}(\ln c, x) = c \cdot x$. Cost: $1n$.

\noindent\textbf{Affected:} $\lambda \cdot \mu$ when both $\lambda,\mu$ are variable probabilities (not forced to $(0,1)^2$
as a domain restriction --- F16 circuit must work on all reals).

\subsection{Selected Expression Re-Audit}

\begin{longtable}{p{4.5cm}cccp{4.5cm}}
\toprule
Expression & Old F16 & New F16 & $\Delta$ & Notes \\
\midrule
\endhead
\midrule
\multicolumn{5}{r}{\textit{continued}} \\
\endfoot
\bottomrule
\endlastfoot
\multicolumn{5}{l}{\textbf{--- Unchanged (no free-variable mul/div) ---}} \\
$e^x$ & $1n$ & $1n$ & $0$ & F16 primitive \\
$\ln x$ & $1n$ & $1n$ & $0$ & F16 primitive \\
$\sqrt{x}$ ($x>0$) & $1n$ & $1n$ & $0$ & EPL(0.5,$x$) \\
$x^n$ ($x>0$) & $1n$ & $1n$ & $0$ & EPL($n$,$x$) \\
$e^{-x}$ & $2n$ & $2n$ & $0$ & neg+exp \\
$e^x - 1$ & $1n$ & $1n$ & $0$ & F16 sign-variant \\
Softplus $\ln(1+e^x)$ & $2n$ & $2n$ & $0$ & B: sign-variant+ln \\
Mayer $f$: $e^{-\beta u} - 1$ & $3n$ & \warn{$3n$--$4n$} & $?$ & $\beta$ const: $3n$; $\beta$ free: $4n$ \\
$e^{x-y}$ & $1n$ & $1n$ & $0$ & $D_{F1}(x,y)$ directly \\
$\ln(x/y)$ ($x,y>0$) & $3n$ & $3n$ & $0$ & sub+ln = 2n, no div needed \\
$\ln(xy)$ ($x,y>0$) & $3n$ & $3n$ & $0$ & add+ln \\
$x \cdot e^y$ & $3n$ & $3n$ & $0$ & $D_{F14}(y,x)$ is $1n$!  Actually $e^y \cdot x$? Let me check: $D_{F14}(y,x) = \exp(y + \ln x) = e^y \cdot x$. Yes, $1n$. \\
\multicolumn{5}{l}{\textbf{--- Affected (free-variable multiplication) ---}} \\
$\lambda \ln(\lambda/\mu)$ (KL term) & $5n$ & $5n$ & $0$ & Route: sub(ln)+mul [2+3]; same via alg.\ identity \\
$\lambda \mu$ (two free real vars) & $2n$ & $3n$ & $+1$ & Verified: no 2-node F16 \\
$(p-q)\ln(p/q)$ (Jeffreys) & $7n$ & $8n$ & $+1$ & Extra mul by free var \\
$-\lambda\ln\lambda$ (vN entropy) & $4n$ & $4n$ & $0$ & $\ln$ then mul by $\lambda$: 1+3=4; same \\
$\sqrt{\lambda\mu}$ (Bures pos.) & $3n$ & $3n$ & $0$ & mul pos=1n + EPL = 2n total; wait: $D_{F16}(\lambda,\mu)=\lambda\mu$[1n], EPL[1n] = $2n$ \\
\multicolumn{5}{l}{\textbf{--- Affected (general-domain division) ---}} \\
$x / y$ (general real) & $2n$ & $\geq 3n$ & $+1$ & Lean-proved \\
$e^x / e^y$ (general) & $3n$ & $3n$ & $0$ & Sub-then-exp avoids div \\
Softmax $e^x/(e^x+e^y)$ (general) & $6n$ & $\geq 7n$ & $+1$ & Requires general div \\
$\lambda/\mu$ (KL: ratio step) & $2n$ & $\geq 3n$ & $+1$ & But ln-of-ratio route avoids this \\
\multicolumn{5}{l}{\textbf{--- Category C (notation-only, unchanged at F16 baseline) ---}} \\
$e^x + e^y$ (EEA) & $4n$ & $4n$ & $0$ & No F16 shortcut \\
LSE $\ln(e^x+e^y)$ & $5n$ & $5n$ & $0$ & No F16 shortcut \\
$\cosh x$ & $6n$ & $6n$ & $0$ & No F16 shortcut \\
$\sinh x$ & $6n$ & $6n$ & $0$ & No F16 shortcut \\
\end{longtable}

\section{KL Divergence: A Key Case Study}

The KL term $\lambda \ln(\lambda/\mu)$ appeared in the previous audit as $5n$.
The route was: $\mathrm{div}(\lambda,\mu)[2n] + \ln[1n] + \mathrm{mul}(\lambda,\cdot)[2n] = 5n$.
With corrected bounds: $\mathrm{div}[3n] + \ln[1n] + \mathrm{mul}[3n] = 7n$.

However, the algebraic identity $\ln(\lambda/\mu) = \ln\lambda - \ln\mu$ avoids the division:
\[
\lambda\ln(\lambda/\mu) = \lambda(\ln\lambda - \ln\mu) = \lambda\cdot\mathrm{sub}(\ln\lambda,\ln\mu).
\]
Route: $\ln\lambda[1n] + \ln\mu[1n] + \mathrm{sub}[2n] + \mathrm{mul}(\lambda,\cdot)[3n] = 7n$.
Both routes give $7n$. Previous claim of $5n$ is not achievable without $\mathrm{mul}=2n$.

\begin{finding}
KL divergence per term: $5n \to 7n$ under corrected bounds.
\end{finding}

\section{Boltzmann Ratio: A Case Where Nothing Changes}

$e^{-\beta(E_j - E_i)} = e^{\beta(E_i - E_j)}$. If $\beta$ is a fixed parameter:
\[
\mathrm{sub}(E_i, E_j)[2n] \to \text{mul}(\beta, \cdot): D_{F14}(\ln\beta, \cdot)[1n] \to \exp[0n, absorbed].
\]
Total: $3n$ for the ratio. No change, because $\beta$ is a constant, not a free variable.

\begin{finding}
Boltzmann ratio with fixed $\beta$: unchanged at $3n$ F16 (or $6n$ for both factors, algebraic save).
\end{finding}

\section{New Potential Collapses}

With the bounds now locked, we look for new reductions. A ``collapse'' would be an expression
previously thought to require $k$ nodes that can be achieved in fewer.

\subsection{div-then-ln avoidance}

Any expression of the form $\ln(x/y)$ can avoid the general-domain division:
$\ln(x/y) = \ln x - \ln y$, costing $1n + 1n + 2n = 4n$ via sub, versus $\mathrm{div}[3n] + \ln[1n] = 4n$. Same cost, but the sub route avoids the expensive general division node and has no domain restriction.

\subsection{mul-then-exp identities}

Expressions of the form $x \cdot e^y$ can be captured as a single F16 node:
$D_{F14}(y, x) = \exp(y + \ln x) = e^y \cdot x$ for $x > 0$. This is $1n$ total --- a genuine collapse
not visible in the old audit's "mul+exp = 3n" counting.

For general (possibly negative) $x$: $x \cdot e^y = \mathrm{sign}(x) \cdot |x| \cdot e^y$,
which requires extra nodes, confirming the $3n$ lower bound is tight for fully general inputs.

\subsection{Power-of-ratio expressions}

$(x/y)^n$ for $x,y > 0$: $D_{F13}(n, D_{F16}(x, D_{F13}(-1,y))) = (x/y)^n$. Cost: $3n$.
Previously might have been counted as div($2n$) + pow($1n$) = $3n$ --- same, no change.

\section{Summary of Impact}

\begin{center}
\begin{tabular}{lrr}
\toprule
Category & Expressions affected & Direction \\
\midrule
Free-variable mul (general domain) & All with $x \cdot y$, $x,y$ free & $+1n$ \\
General-domain division & All with $x/y$, $x,y$ free & $+1n$ \\
Positive-domain division & $1n \to 2n$ (v5.2 error corrected) & $+1n$ \\
Constant-factor expressions & $c \cdot x$, $c$ fixed & No change (EPL) \\
Sub-then-exp, sub-then-ln routes & No free mul/div & No change \\
Category C (EEA, LLA etc.) & All notation-only & No change \\
\midrule
\textbf{Core 15n positive total} & & \ok{Preserved} \\
\textbf{General domain total} & Recount to $\approx 18$--$20n$ & \bad{Up from ``22n''} \\
\bottomrule
\end{tabular}
\end{center}

\noindent\textbf{Paradox note}: The general-domain total appears to \emph{decrease} from 22n to
$\approx 18n$ despite mul and div increasing. This is because the original ``22n'' figure likely
included incorrect overestimates elsewhere, or counted more operations. A precise recount
is performed in Core10\_ReCheck.tex.

\end{document}
